\chapter{Station Backtest}

\section{Mandat}
La station Backtest evalue une strategie complete dans un cadre
\textit{Walk-Forward}. Son objectif est de mesurer le transfert de performance
entre in-sample et out-of-sample, et non de maximiser une performance historique
ponctuelle \cite{pardo2008}.

\section{Pipeline Walk-Forward}
Le cycle de validation est compose des etapes suivantes :
\begin{enumerate}
\item generation de variantes parametriques faisables,
\item optimisation sur fenetre IS,
\item test du gagnant sur fenetre OOS,
\item repetition par roulage temporel,
\item selection de la configuration la plus robuste.
\end{enumerate}

\section{Fonctions de decision}
\subsection{Walk-Forward Efficiency}
Le critere de transfert principal est :
\begin{equation}
\text{WFE} = \frac{R_{\mathrm{OOS}}}{R_{\mathrm{IS}}} \times 100
\label{eq:wfe}
\end{equation}
avec $R_{\mathrm{OOS}}$ et $R_{\mathrm{IS}}$ les rendements annualises respectifs.
L'equation~\ref{eq:wfe} est utilisee comme gate de robustesse (\(\text{WFE}\geq 50\%\))
\cite{pardo2008}.

\subsection{Sizing issu des transactions OOS}
Le sizing final repose sur le critere de Kelly :
\begin{equation}
f^\star = p - \frac{1-p}{R}
\label{eq:kelly}
\end{equation}
ou \(p\) est le taux de gain et \(R\) le ratio gain/perte moyen, estimes sur les
trades OOS concatenes \cite{kelly1956,thorp2006}. L'equation~\ref{eq:kelly}
evite d'injecter des hypothese subjectives de sizing.

\section{Validation statistique}
Au-dela des metriques de performance, la station integre des controles de
significativite (tests de permutation / Monte Carlo) et des corrections de biais
de selection de type DSR \cite{bailey2014}.

\section{Criteres d'acceptation}
Le Tableau~\ref{tab:acceptance-backtest} resume les criteres de decision.

\begin{table}[H]
\centering
\caption{Criteres d'acceptation/rejet d'une strategie}
\label{tab:acceptance-backtest}
\begin{tabular}{p{4.5cm}p{8.8cm}}
\toprule
\textbf{Critere} & \textbf{Interpretation} \\
\midrule
WFE & Doit depasser un seuil minimal de transfert IS $\rightarrow$ OOS. \\
Rentabilite OOS majoritaire & La majorite des fenetres OOS doit rester profitable. \\
Stabilite parametrique & Le profil d'optimisation doit presenter des plateaux, pas des pics isoles. \\
Significativite statistique & p-value et DSR doivent rester compatibles avec une hypothese de robustesse. \\
\bottomrule
\end{tabular}
\end{table}

Le Tableau~\ref{tab:acceptance-backtest} formalise la regle metier ``pas de
production sans robustesse''.

